\section{Introduction and main theorem}\label{sec:introduction}

Frobenius-twisted fixed points on a nonproper variety need not agree with
the compactly supported cohomological trace at small Frobenius powers.
For a nonlinear polynomial automorphism of the affine plane, the discrepancy
can be measured by intersecting its graph with the Frobenius graph at
infinity. We carry out that measurement for every iterate of a generalized
H\'enon map and every nonresonant Frobenius power.

Fix a prime power $q=p^e$ and put $k=\overline{\F}_q$. Throughout the paper,
\begin{equation}\label{eq:map}
 H:\A^2_k\longrightarrow\A^2_k,
 \qquad H(x,y)=(y,f(y)-ax),
 \qquad f\in\F_q[y],\quad \deg f=d\geq2,\quad a\in\F_q^*.
\end{equation}
The inverse is the polynomial map
\begin{equation}\label{eq:inverse}
 H^{-1}(x,y)=\bigl((f(x)-y)/a,x\bigr).
\end{equation}
There is no monicity assumption and no restriction on a lower coefficient
of $f$. For a positive integer $r$, let
\begin{equation}\label{eq:frobenius}
 \Phi_r(x,y)=(x^{q^r},y^{q^r}).
\end{equation}
This is the actual power morphism over $k$, obtained by base change from
the $q^r$-Frobenius morphism over $\F_q$. We use the explicit convention
\eqref{eq:frobenius} throughout, independently of terminological conventions
for arithmetic and geometric Galois Frobenius.

For positive integers $n$ and $r$, define
\begin{equation}\label{eq:equalizer}
 \begin{aligned}
 E_H(n,r)&=\{P\in\A^2_k:H^n(P)=\Phi_r(P)\}
       \quad\text{as an equalizer scheme},\\
 M_H(n,r)&=\#E_H(n,r)(k).
 \end{aligned}
\end{equation}
The two integers have different roles: $n$ is the H\'enon iteration clock
and $r$ is the Frobenius clock. In particular, \eqref{eq:equalizer} is not
the count $\#\Fix(H^n)(\F_{q^r})$ of rational points on one fixed-period
scheme. Since $H$ is invertible, it is also the geometric fixed-point scheme
of $H^{-n}\circ\Phi_r$.

\begin{theorem}[Exact nonresonant equalizer count]\label{thm:main}
Let $H$ satisfy \eqref{eq:map}. For every $n,r\geq1$, the affine equalizer
$E_H(n,r)$ is finite and reduced. Set $D=d^n$ and $Q=q^r$. If $D\ne Q$,
then
\begin{equation}\label{eq:maxlaw}
 \boxed{M_H(n,r)=\max\{DQ,Q^2\}.}
\end{equation}
More precisely, the closures of the graphs of $H^n$ and $\Phi_r$ in
$\Pj^2_k\times\Pj^2_k$ meet properly. Their only intersections outside
$\A^2_k\times\A^2_k$ are
$(\Imn,\Imn)$ and $(\Ip,\Ip)$, where
$\Imn=[0:1:0]$ and $\Ip=[1:0:0]$. Their respective local intersection
multiplicities are $Q\min(D,Q)$ and $1$.
\end{theorem}

\begin{corollary}[All pairs of clocks]\label{cor:allclocks}
If $d$ is not a power of $p$, then \eqref{eq:maxlaw} holds for every
$n,r\geq1$ and every map \eqref{eq:map} of degree $d$ over $\F_q$.
\end{corollary}

The theorem supplies two connected conclusions. First, it determines the
complete coefficient-uniform boundary contribution, including the regime
$D>Q$ in which the affine count exceeds $Q^2$. Second, it turns that
contribution into an exact threshold and distinguishes fixed-time
Frobenius slices from the iterates of one fixed map. The equal-degree
case $D=Q$ is excluded from the counting formula itself, not just from
the proof's presentation.

The total projective intersection is $1+DQ+Q^2$. The substantive local
calculation is therefore the length $Q\min(D,Q)$ at the infinity point
where $H^n$ is regular. At the opposite infinity point the inverse graph
is regular, and the Frobenius differential forces length one. This
asymmetry explains why using only the largest affine degrees would miss
the correct count.

Section~\ref{sec:context} states the classical ownership of the trace
principles used for interpretation. Sections~\ref{sec:geometry}
and~\ref{sec:boundary} prove the theorem and corollary.
Section~\ref{sec:slices} derives the exact threshold, the fixed-time
slice formula and a finite-dimensional trace obstruction.
Section~\ref{sec:diagonal} treats genuine diagonal dynamics and the
resonant counterexample. All proofs of the stated results are in the
main text.
